\documentclass[11pt]{article}
\usepackage[margin=0.82in]{geometry}
\usepackage[T1]{fontenc}
\usepackage{lmodern}
\usepackage{amsmath,amssymb,mathtools}
\usepackage{graphicx}
\usepackage{xcolor}
\usepackage{booktabs}
\usepackage{tabularx}
\usepackage{microtype}
\usepackage{hyperref}
\usepackage{caption}
\usepackage{fancyhdr}
\usepackage{enumitem}
\usepackage{array}
\graphicspath{{figures/}}
\hypersetup{colorlinks=true,linkcolor=black,urlcolor=blue,citecolor=black}
\setlength{\parindent}{0pt}
\setlength{\parskip}{0.58em}
\setlist{nosep,leftmargin=1.5em}
\pagestyle{fancy}
\setlength{\headheight}{14pt}
\fancyhf{}
\lhead{Jacobian Counter-example and Elliptic Integrals}
\rhead{July 2026}
\cfoot{\thepage}
\newcommand{\dd}{\,d}
\newcommand{\RR}{\mathbb R}
\definecolor{softblue}{RGB}{233,242,250}
\definecolor{softgold}{RGB}{250,245,228}
\definecolor{softred}{RGB}{252,238,238}
\newcommand{\bluebox}[1]{\begin{center}\colorbox{softblue}{\parbox{0.93\linewidth}{#1}}\end{center}}
\newcommand{\goldbox}[1]{\begin{center}\colorbox{softgold}{\parbox{0.93\linewidth}{#1}}\end{center}}
\newcommand{\redbox}[1]{\begin{center}\colorbox{softred}{\parbox{0.93\linewidth}{#1}}\end{center}}

\title{\vspace{-1.2em}\textbf{Jacobian Counter-example and Elliptic Integrals}}
\author{Harm.On.ica S-O-L 5.6 (Open AI), Bradley Klee}
\date{July 2026}

\begin{document}
\maketitle
\vspace{-1.8em}

\begin{center}
\includegraphics[width=0.99\textwidth,keepaspectratio]{nested_curves_headliner_clean.png}
\end{center}
\vspace{-0.8em}

\begin{abstract}
A two-dimensional slice of the 2026 counterexample to the Jacobian conjecture
starts with nested closed curves inside a triangle and maps the filled regions
into three-dimensional space. Two areas must be kept separate. The
\emph{pre-image area} is the flat area before deformation. The \emph{image
area} is the true area of the bent and stretched sheet afterward. The
pre-image area has a direct hypergeometric formula and leads, through OEIS
A007004 and its derivative A006480, to Ramanujan's cubic period function. For
the image area we give the full deformation density, an exact two-dimensional
integral, and an exact one-dimensional integral for its derivative. Expanding
the deformation density near the center produces a convergent sequence of
finite linear combinations of hypergeometric functions. The first twenty area
coefficients, direct numerical integration, and ordinary triangle meshes all
agree. The final section poses a geometric question: can the flat triangle be
folded and stretched continuously until its three corners meet, and can the
given polynomial map be reached as a limit of small smooth deformations?
\end{abstract}

\clearpage
\section{From the Jacobian conjecture to the triangle map}

The Jacobian determinant records how a map changes a tiny area or volume. If
it never vanishes, the map can be reversed near every point. The Jacobian
conjecture asked whether a polynomial map with one fixed nonzero Jacobian
determinant must also have one global polynomial inverse. In July 2026,
Levent Alp\"oge announced a three-variable counterexample: its Jacobian
determinant is constantly $-2$, yet three different inputs have the same
output~\cite{alpoge,speyer}.

The map is
\[
F(x,y,z)=(P,Q,R),
\]
with
\begin{align*}
P&=(1+xy)^3z+y^2(1+xy)(4+3xy),\\
Q&=y+3x(1+xy)^2z+3xy^2(4+3xy),\\
R&=2x-3x^2y-x^3z.
\end{align*}
The three colliding inputs lie in the plane
\[
\iota(u,v)=\left(u-v,-\frac32(u-v),\frac{27}{4}(u+v)-\frac14\right).
\]
Thus the two-variable map studied here is $G=F\circ\iota$. Its triangle
corners satisfy
\[
G(0,0)=G(1,0)=G(0,1)=\left(-\frac14,0,0\right).
\]
Filling the triangle and applying $G$ turns the logical counterexample into a
visible curved sheet. The nested curves below are level curves in that
triangle and their images under $G$.

\section{The picture before the integrals}

Start with the triangle whose corners are
\[
(0,0),\qquad(1,0),\qquad(0,1)
\]
in coordinates $(u,v)$. Inside it, consider the function
\[
H(u,v)=uv(1-u-v).
\]
Its largest value is $1/27$ at the center $(1/3,1/3)$. We label the nested
closed loops by
\[
H(u,v)=\frac{1-m}{27},\qquad 0\leq m\leq1.
\]
At $m=0$ the loop has shrunk to the center. As $m$ grows, it expands. At
$m=1$ it reaches the three sides of the triangle. We call the filled region
inside the loop $D_m$ and its boundary $C_m$.

A polynomial map $G$ sends the starting plane into three-dimensional space:
\[
G(u,v)=(U(u,v),V(u,v),W(u,v)).
\]
The important geometric fact is that the three triangle corners have the same image,
\[
G(0,0)=G(1,0)=G(0,1)=\left(-\frac14,0,0\right).
\]
Thus the final triangular boundary becomes three loops meeting at one point.
For $m<1$, the closed oval $C_m$ becomes the space curve
\[
\Gamma_m=G(C_m),
\]
and the filled region becomes the mapped sheet
\[
\Sigma_m=G(D_m).
\]
The map used here is the restriction of the explicit polynomial map announced
by Levent Alp\"oge in July 2026~\cite{alpoge}.

\section{Two different meanings of area}

Imagine drawing a closed shape on a flat, stretchable sheet. Measure the inked
region before touching the sheet: this is the \emph{pre-image area}. Now bend
and stretch the sheet according to the map and measure the same marked region
along the deformed surface: this is the \emph{image area}. The second number
includes the stretching caused by the map; it is not merely a different view
of the first.

\begin{center}
\renewcommand{\arraystretch}{1.25}
\begin{tabularx}{0.98\textwidth}{>{\bfseries}p{0.18\textwidth}X X}
\toprule
Area & When it is measured & What it tells us\\
\midrule
Pre-image area
& Before applying the deformation map.
& The flat area we started with; it preserves the original timing.\\
Image area
& After deformation, measured directly along the bent and stretched sheet.
& The true area after stretching; it depends on the map but not on how the finished surface is oriented.\\
\bottomrule
\end{tabularx}
\end{center}

The image area is the target of this paper. Changing the deformation map can
change it, because a different map stretches the sheet differently. Rotating
or reflecting the completed surface cannot change it. This combination---map
dependent but orientation invariant---is exactly what we need from a quantity
that belongs to the image.

\section{Integral I: pre-image area}

The pre-image area is the area of the filled loop before the map is applied:
\[
\mathcal A_0(m)=\iint_{D_m}\dd u\,\dd v.
\]
Walking once around the boundary gives the same number:
\[
\mathcal A_0(m)=-\oint_{C_m}v\,\dd u,
\]
with the sign fixed by the direction of travel.

In mechanics, enclosed area divided by $2\pi$ is called an \emph{action}. Here
\[
J_0(m)=\frac{\mathcal A_0(m)}{2\pi}.
\]
The exact answer is especially simple after normalization. Define
\[
F_3(m)={}_2F_1\!\left(\frac13,\frac23;1;m\right)
=\sum_{n\geq0}\frac{(3n)!}{(n!)^3}\left(\frac{m}{27}\right)^n
\]
and its integral
\[
S_3(m)=\int_0^m F_3(t)\,\dd t.
\]
Then
\[
\boxed{
\mathcal A_0(m)=\frac{2\pi\sqrt3}{27}S_3(m),
\qquad
J_0(m)=\frac{\sqrt3}{27}S_3(m).}
\]
This is the exact quantity that preserves the original clock.

The cubic rescaling has integer coefficients:
\[
\frac{S_3(27x)}{27}
=x+3x^2+30x^3+420x^4+6930x^5+126126x^6+\cdots.
\]
The coefficient of $x^{n+1}$ is
\[
\frac{(3n)!}{(n+1)(n!)^3},
\]
which is OEIS A007004~\cite{oeis7004}.

Its derivative provides the direct link to Ramanujan. Write
\[
G_7(x)=\sum_{n\geq0}\frac{(3n)!}{(n+1)(n!)^3}x^{n+1}
=x\,{}_2F_1\!\left(\frac13,\frac23;2;27x\right).
\]
Then
\[
G_7'(x)=\sum_{n\geq0}\frac{(3n)!}{(n!)^3}x^n
={}_2F_1\!\left(\frac13,\frac23;1;27x\right).
\]
The derivative coefficients are OEIS A006480. OEIS identifies this
hypergeometric function with period functions for elliptic-curve families and
connects it to Ramanujan's work on modular equations and approximations to
$\pi$~\cite{oeis6480,ramanujan}. Thus A007004 is an integrated version of the
Ramanujan period sequence A006480.

\goldbox{The pre-image area is exact and preserves the original clock, but it
does not measure how much the map stretches the sheet. It is the area before
deformation, recovered from the image by undoing that deformation.}

\subsection*{How a triangle mesh recovers the pre-image area}

The mesh is a later numerical check, not the definition of the integral. A
small mapped triangle with vertices $\mathbf r_1,\mathbf r_2,\mathbf r_3$ has
an area arrow
\[
\mathbf a_\triangle=
\frac12(\mathbf r_2-\mathbf r_1)\times(\mathbf r_3-\mathbf r_1).
\]
Its length is the triangle's image area, while its direction records how the
triangle is facing. To recover the pre-image area instead, use the correcting
vector
\[
B_0(u,v)=
\frac{G_u\times G_v}{\|G_u\times G_v\|^2},
\qquad
B_0\cdot(G_u\times G_v)=1.
\]
Then
\[
\sum_{\triangle} B_0(u_\triangle,v_\triangle)\cdot
\mathbf a_\triangle
\longrightarrow \mathcal A_0(m).
\]
These weights undo the local stretching of the map. They recover the area
before deformation rather than measuring the size of the deformed sheet.

\section{Integral II: image area after stretching}

The image area is the ordinary area of the mapped sheet after the map has bent
and stretched it:
\[
\boxed{
\mathcal A_\Sigma(m)=
\iint_{D_m}\|G_u\times G_v\|\,\dd u\,\dd v.}
\]
This number is unchanged if the completed surface is rotated or reflected,
but it changes when the deformation map changes. It therefore measures the
geometry created by the map itself.

For an exact formula, use the centered coordinates
\[
u=\frac13+\frac{x}{3}+\frac{y}{3\sqrt3},
\qquad
v=\frac13-\frac{x}{3}+\frac{y}{3\sqrt3}.
\]
The filled region is
\[
D_m=\{(x,y):M(x,y)\le m\},
\qquad
M=x^2+y^2+\frac{2}{3\sqrt3}y(y^2-3x^2).
\]
Write
\[
\mathcal J(x,y)=\|G_x\times G_y\|
=\sqrt{C_1(x,y)^2+C_2(x,y)^2+C_3(x,y)^2},
\]
where
\begin{align*}
C_1={}&-\frac{4\sqrt3x}{81}
\bigl(24\sqrt3x^6y-4x^6-72x^4-54\sqrt3x^2y-27x^2+81\bigr),\\
C_2={}&-\frac{\sqrt3}{81}
\bigl(48\sqrt3x^8y-8x^8-72\sqrt3x^6y-132x^6
-108\sqrt3x^4y+162x^4\\
&\hspace{3.1cm}+162\sqrt3x^2y+297x^2-162\bigr),\\
C_3={}&-\frac{\sqrt3(2x^2-3)^2}{324}
\bigl(48\sqrt3x^6y-8x^6-72\sqrt3x^4y-132x^4
-108\sqrt3x^2y\\
&\hspace{3.1cm}+162x^2+162\sqrt3y+405\bigr).
\end{align*}
Thus the deformed area has the explicit exact form
\[
\boxed{
\mathcal A_\Sigma(m)=
\iint_{M(x,y)\le m}\mathcal J(x,y)\,\dd x\,\dd y.}
\]
Nothing is hidden in the norm notation: the three polynomials above are the
three components of the local area-stretching vector.

There is also an exact one-dimensional formula for the rate at which area is
added. In polar coordinates,
\[
M(r,\theta)=r^2-\frac{2}{3\sqrt3}r^3\sin(3\theta).
\]
Let $R_m(\theta)$ be the positive solution of
\[
R_m^2-\frac{2}{3\sqrt3}R_m^3\sin(3\theta)=m.
\]
Differentiating the radial integral gives
\[
\boxed{
\mathcal A_\Sigma'(m)=\frac12\int_0^{2\pi}
\frac{\mathcal J(R_m\cos\theta,R_m\sin\theta)}
{1-\dfrac{R_m}{\sqrt3}\sin(3\theta)}\,\dd\theta.}
\]
The double integral gives the accumulated image area; this single integral
gives the new area contributed by the next thin shell.

At the center,
\[
J_c=\mathcal J(0,0)=\frac{\sqrt{6267}}4.
\]
Normalize by
\[
\Psi(m)=\frac{\mathcal A_\Sigma(m)}{\pi J_c}.
\]
Then
\[
\Psi(m)=\frac{1}{\pi J_c}\int_0^{2\pi}
\int_0^{R_m(\theta)}\mathcal J(r\cos\theta,r\sin\theta)
\,r\,\dd r\,\dd\theta,
\]
and
\[
\Psi'(m)=\frac{1}{2\pi J_c}\int_0^{2\pi}
\frac{\mathcal J(R_m\cos\theta,R_m\sin\theta)}
{1-\dfrac{R_m}{\sqrt3}\sin(3\theta)}\,\dd\theta.
\]
These are exact formulas for the deformed area and its derivative.

\subsection*{Hypergeometric form near the center}

The same cubic hypergeometric functions that describe the pre-image area also
organize the local expansion of the image area. Put
\[
F_3(m)={}_2F_1\!\left(\frac13,\frac23;1;m\right),
\qquad
S_3(m)=m\,{}_2F_1\!\left(\frac13,\frac23;2;m\right).
\]
Expand the exact stretching density at the center,
\[
\frac{\mathcal J(x,y)}{J_c}=P_0(x,y)+P_1(x,y)+P_2(x,y)+\cdots,
\]
where $P_d$ is the homogeneous part of degree $d$. For each $N$, integrate the
finite polynomial $P_0+\cdots+P_{2N-2}$ over $D_m$. Exact reduction on the
cubic curves gives polynomials $A_N(m)$ and $B_N(m)$ such that
\[
\boxed{
\Psi_N(m)=A_N(m)F_3(m)+B_N(m)S_3(m),
\qquad
\Psi(m)-\Psi_N(m)=O(m^{N+1}).}
\]
As $N$ increases, these finite hypergeometric combinations converge to the
exact image-area function near the center. Differentiation remains in the
same explicit family. If
\[
F_3^+(m)={}_2F_1\!\left(\frac43,\frac53;2;m\right),
\qquad F_3'(m)=\frac29F_3^+(m),
\]
then
\[
\boxed{
\Psi_N'(m)=
\bigl(A_N'(m)+B_N(m)\bigr)F_3(m)
+\frac29A_N(m)F_3^+(m)+B_N'(m)S_3(m).}
\]
The companion file gives the exact $A_8$ and $B_8$, whose combination agrees
with the first eight nonzero coefficients of $\Psi$. This is a practical
hypergeometric form without printing several lines of very large rational
numbers in the main text.

The local series begins
\[
\begin{aligned}
\Psi(m)={}&m
-\frac{8798863}{78550578}m^2
-\frac{656984015341439}{4627644978100563}m^3\\
&-\frac{5888704400396591850310}
{545256281712894848663121}m^4
+\frac{9229642725727364600800502643893}
{856603921733574408314213636678760}m^5+O(m^6).
\end{aligned}
\]
The image-area action is therefore
\[
J_\Sigma(m)=\frac{\mathcal A_\Sigma(m)}{2\pi}
=\frac{\sqrt{6267}}8\Psi(m).
\]

The derivative has a notable integer rescaling. With
\[
C=235651734=54\cdot2089^2,
\]
the first twenty exactly computed coefficients of $\Psi'(Cx)$ are integers.
The first eight are
\[
\begin{aligned}
\Psi'(Cx)={}&1-52793178x-23651424552291804x^2\\
&-565315622438072817629760x^3
+166133569063092562814409047590074x^4\\
&+19584594743364404616475143307241406057420x^5\\
&+4509774622501256294260247296713228343209292240728x^6\\
&+1755635498235399950242562905014394392824488851953334778752x^7
+O(x^8).
\end{aligned}
\]
All twenty integers are included in the companion table.

\subsection*{Independent numerical check}

The exact series, direct evaluation of the integral, and ordinary mapped
triangle areas agree:
\begin{center}
\begin{tabular}{rrrr}
\toprule
$m$ & direct integral & 20-term series & 12,640-triangle mesh\\
\midrule
0.10 & 6.139023 & 6.139023 & 6.137024\\
0.40 & 23.182660 & 23.182660 & 23.169230\\
0.70 & 37.237531 & 37.243141 & 37.204348\\
\bottomrule
\end{tabular}
\end{center}

\begin{figure}[ht]
\centering
\includegraphics[width=0.72\textwidth]{true_surface_mesh_convergence.png}
\caption{Ordinary triangle-area sums converge to direct numerical integration
of the image area at $m=0.70$.}
\label{fig:surface-convergence}
\end{figure}

\section{Folding the triangle continuously}

The last question is geometric rather than computational. Begin with the flat
triangle in the plane and imagine lifting and stretching it continuously until
the three corners meet. The given polynomial map shows the final arrangement:
three distinct corners arrive at one point while each corner still has a
well-defined local sheet. What we would like is a continuous movie of how that
collision can be made.

One precise version is to seek maps
\[
G_t:D_1\longrightarrow\mathbb R^3,
\qquad 0\le t\le1,
\]
with
\[
G_0(u,v)=(u,v,0),
\qquad
G_1(u,v)=G(u,v),
\]
after a harmless change of coordinates in the starting plane. For $t<1$, the
three corners should remain distinct and the sheet should stay smooth. At the
last moment, the corners may meet exactly as they do under $G$.

The equilateral triangle suggests a simple picture. Lift the three corners in
three different directions, as though the triangle were being folded toward
a tetrahedral point, and let those directions move along curves on a sphere.
The sheet is allowed to stretch, so this need not be a rigid paper-folding
problem. The sphere records directions; the triangle records which points are
being joined.

There are two closely related questions for colleagues:
\begin{enumerate}
\item Can one write down such a smooth family $G_t$ directly?
\item Can the final map be obtained as a limit of many small, smooth changes
that each stay close to the identity?
\end{enumerate}
A product of ordinary constant matrices cannot do this by itself: it remains a
single linear map. The useful replacement is a composition of
\emph{position-dependent} near-identity maps, for example
\[
X\longmapsto X+\varepsilon V_k(X),
\]
or the continuous flow
\[
\frac{dX}{dt}=V_t(X).
\]
Their local derivative matrices multiply along each path, while the maps as a
whole can bend and stretch the triangle nonlinearly. An ambient deformation
that is invertible at every stage cannot merge three points, so the collision
must occur in the final limit or in one final non-invertible step. That is not
a defect; it identifies the exact event the construction must explain.

\begin{figure}[ht]
\centering
\includegraphics[width=0.72\textwidth]{mapped_separatrix_3d.png}
\caption{The three boundary arcs meet at one image point. The proposed
continuous folding problem asks for a smooth path from the flat triangle to
this final arrangement.}
\label{fig:critical-arcs}
\end{figure}

\section{Files, checks, and questions for the next step}

The companion files include the exact map, the three polynomial components
$C_1,C_2,C_3$, the first twenty image-area coefficients, the first twenty
rescaled derivative integers, the $N=8$ hypergeometric reduction, direct
integration values, and mesh checks. The main files are
\begin{itemize}
\item \path{data/intrinsic_area_density.txt};
\item \path{data/true_surface_area_series.csv};
\item \path{data/image_area_hypergeometric_reduction_N8.txt};
\item \path{data/true_surface_area_validation.csv};
\item \path{scripts/derive_image_area_hypergeometric.py}.
\end{itemize}

The calculations now provide an exact deformation integral, an exact shell
integral, a convergent hypergeometric reduction near the center, and three
independent numerical checks. The natural next questions are whether the
coefficient functions $A_N,B_N$ approach a simpler named form, whether the
integer rescaling continues to all orders, and whether a smooth family $G_t$
can make the triangle's three corners meet in the final limit.

\begin{thebibliography}{99}
\bibitem{alpoge}
Levent Alp\"oge,
``hello there the jacobian conjecture is false ...'', public announcement,
July 19, 2026,
\url{https://x.com/__alpoge__/status/2079028340955197566}.

\bibitem{crawford}
Tom Crawford,
``Oxford Calculus: Jacobians Explained,''
YouTube video,
\url{https://www.youtube.com/watch?v=YqMelRryG8U}.

\bibitem{mathgeek}
Mathgeek,
``The Jacobian Conjecture,''
YouTube video, July 2, 2023,
\url{https://www.youtube.com/watch?v=Jxu3is_hrZI}.

\bibitem{origin}
L\'azaro O. Rodr\'iguez D\'iaz,
``On the origin of the Jacobian conjecture,''
2025,
\url{https://arxiv.org/abs/2512.23614}.

\bibitem{countervideo}
``Counter Example To Jacobian Conjecture Found By Levent Alpoge,''
YouTube video, July 2026,
\url{https://www.youtube.com/watch?v=2ausgtmjn6g}.

\bibitem{detailsvideo}
``AI kills the Jacobian conjecture -- the details,''
YouTube video, July 2026,
\url{https://www.youtube.com/watch?v=8OBxOuyDBU8}.

\bibitem{speyer}
David Speyer,
``The new counterexample to the Jacobian conjecture,''
Secret Blogging Seminar, July 20, 2026,
\url{https://sbseminar.wordpress.com/2026/07/20/the-new-counterexample-to-the-jacobian-conjecture/}.

\bibitem{oeis7004}
N. J. A. Sloane and James Propp,
``A007004: $(3n)!/((n+1)(n!)^3)$,''
The On-Line Encyclopedia of Integer Sequences,
\url{https://oeis.org/A007004}.

\bibitem{oeis6480}
N. J. A. Sloane,
``A006480: $(3n)!/(n!)^3$,''
The On-Line Encyclopedia of Integer Sequences,
\url{https://oeis.org/A006480}.

\bibitem{ramanujan}
Srinivasa Ramanujan,
``Modular Equations and Approximations to $\pi$,''
\emph{Quarterly Journal of Mathematics}, Vol. 45 (1914), 350--372.
\end{thebibliography}

\end{document}
